% -*- mode:LaTex; mode:visual-line; mode:flyspell; fill-column:75-*-

\chapter{Conclusions} \label{chapConclusions}

This thesis has addressed key challenges in realizing efficient bio-inspired locomotion on physical aquatic robots. Although prior work has produced a wide range of biomimetic swimming platforms, the joint design of morphology, compliance, dynamics, and onboard sensing has often remained fragmented. This thesis has argued that aquatic locomotion is best understood as a problem of embodied multiphysics, in which the robot body, compliant structures, actuators, sensors, and surrounding fluid form a coupled dynamical system. From this perspective, meaningful progress requires not only new mechanisms or controllers, but also models and experimental platforms that explicitly capture the interactions among body design, fluid loading, actuation, and sensing. Building on this view, the thesis developed contributions across programmable compliance, variational modeling, optimization-based locomotion, mechatronic platform design, and distributed onboard sensing.

The central contribution of the thesis is the programmable online stiffness modulation (POSM) framework, which converts joint stiffness from a fixed mechanical property into a within-cycle control variable. A modular multi-joint mechanism was developed in which sliding-laminate flexures provide time-varying, joint-specific compliance through simple two-channel actuation. This mechanism was paired with a variational-mechanics model incorporating quasi-steady hydrodynamic forces, enabling gradient-based trajectory optimization over both appendage motion and joint-stiffness schedules. Hardware experiments on a fixed-base testbed and on an untethered fish-inspired robot capable of both body--caudal-fin (BCF) undulation and drag-based pectoral-fin paddling demonstrated that programmable joint-specific stiffness can improve propulsive performance over static and spatially uniform baselines, achieving up to 89\% higher thrust in drag-based swimming, up to 19\% higher mean thrust in BCF swimming, and up to 34.8\% reduction in cost of transport.

Beyond reduced-order hydrodynamic modeling, this thesis also extended the variational modeling perspective toward strongly coupled fluid--robot multiphysics. A unified, differentiable simulator was developed in which the no-slip interface between the robot and the surrounding fluid is treated as an integral constraint, enabling gait optimization for both sustained undulatory swimming and aggressive transient maneuvers such as C-start escapes. To evaluate these modeling and locomotion methods in real water, the thesis introduced two complementary mechatronic platforms: a modular undulatory swimming robot for steady and transient swimming experiments, and SwimGym, a low-cost, open-source drag-based swimming system with an associated digital-twin simulation environment. Together, these platforms provide experimental foundations for studying sim-to-real transfer, model-based gait optimization, and the role of embodied morphology in aquatic locomotion.

The thesis further broadened the embodied-sensing component from articulated aquatic swimmers to modular continuum robots. A reconfigurable continuum platform was developed using 3D-printed compliant joints with finite-element-characterized stiffness, together with a self-contained pose-estimation pipeline that maps onboard magnetic measurements to inter-segment orientations. This system demonstrated that distributed local sensing can reduce dependence on external tracking infrastructure across traversal, grasping, and large-deformation manipulation tasks. In parallel, printable thermoelectric ink was examined as a candidate material technology for thermally responsive robotic devices, extending the embodied-design theme toward soft, multifunctional substrates.

Taken together, the work in this thesis demonstrates that programmable compliance, variational modeling, optimization-based propulsion, embodied platform design, and distributed sensing form a coherent design space for bio-inspired aquatic robots. Rather than treating morphology, dynamics, and sensing as separate modules, the thesis shows how they can be jointly exploited to improve locomotion performance and to support more reliable model-based design. Looking forward, the most natural extensions include closed-loop control of POSM under unsteady inflow and obstacle interaction, integration of strongly coupled fluid--robot multiphysics directly into stiffness optimization, and the use of distributed local flow and proprioceptive sensing to enable fully autonomous bio-inspired swimming beyond controlled laboratory tanks.
